
\subsubsection{Spectral Normalization}

Miyato et al. (2018) introduced spectral normalization for GAN discriminators, constraining weight matrix spectral norms via power iteration. Their method operates on static weight tensors, rescaling W as W_SN = W/σ(W) at each training step. This differs from our approach, which attempts to regularize dynamic Jacobians J_ℓ = ∂h_ℓ/∂h_{ℓ-1} computed via autodiff through attention and layer normalization operations.

\subsubsection{Low-Rank Adaptation}

LoRA (Hu et al., 2021) represents weight updates as low-rank decompositions ΔW = BA with fixed rank r. ARD-LoRA (Shinwari et al., 2025) extends this with automatic rank determination, achieving 99.3\% performance with 0.32\% parameters. These methods operate on small adapter matrices during fine-tuning, not on layer Jacobians during pretraining.

\subsubsection{KV Cache Optimization}

KV-CAT (Gelberg et al., 2026) trains transformers for KV cache compressibility via auxiliary losses on activation statistics. Flash Attention (Dao et al., 2022) achieves 2-4× speedup through IO-aware computation. These approaches optimize specific efficiency dimensions independently.

\subsubsection{Trace Estimation}

Hutchinson's estimator approximates Tr(A) via E[z^T Az] for random vectors z. Bekas et al. (2007) establish O(1/ε²) sample complexity for ε-accuracy, which for coefficient of variation below 15\% at 768 dimensions implies approximately 100+ probes. Our implementation used 10 probes.
